\begin{table}[htbp]\centering
\begin{threeparttable}
\caption{Protein-coding genes on the rDNA species complexes: smallest distance (\%) between the two species, with the
largest distance within either species in brackets.}
\label{tab:pcgenes}
\footnotesize
\begin{tabular}{p{0.44\linewidth}p{0.14\linewidth}p{0.14\linewidth}p{0.14\linewidth}}
\toprule
Pair & Glomalin & RPB1 & H$^+$-ATPase\\
\midrule
\spn{D.\ aestuarii} / \spn{D.\ varaderana} & 0.73 (0.0) & 1.24 (1.73) & 1.22 (0.0)\\
\spn{D.\ arenaria} / \spn{D.\ jakucsiae} & 0.29 (--) & 0.57 (--) & 1.55 (--)\\
\spn{G.\ chinense} / \spn{G.\ rugosae} & 0.73 (0.73) & 0.33 (7.93) & 1.00 (1.32)\\
\spn{R.\ irregularis} / \spn{[R.]\ vesiculiferum} & 1.32 (1.61) & 0.81 (7.54) & 1.18 (5.77)\\
\bottomrule
\end{tabular}
\begin{tablenotes}\footnotesize
\item Sequences of Stefani et al.\ (2025). Indicative species thresholds from that study: 1.0\% / 1.1\% / 1.7\% (glomalin / RPB1 /
H$^+$-ATPase). Distances are edlib infix edit distances including any intron, so they are comparable with each other but
only approximately with those thresholds. Records are joined to species through the reference culture names; records whose
GenBank name disagrees with the reference name of their culture are left out (6).
\end{tablenotes}
\end{threeparttable}
\end{table}
